\begin{abstractpage}

\begin{abstract}{romanian}
Lucrarea descrie o aplicație web pentru antrenamente de forță, construită în jurul unei idei simple:
istoricul de antrenament trebuie să rămână corect în timp. Atunci când un antrenament este efectuat,
planul din care a pornit este copiat („înghețat”) sub formă de instantanee, așa că dacă utilizatorul
modifică sau șterge mai târziu datele de la care a plecat, ce a fost deja înregistrat nu se mai
schimbă. Aplicația permite gestionarea programelor de antrenament, executarea sesiunilor, urmărirea
progresului, un marketplace de programe (publicare, vot, salvare, copiere), o analiză a structurii
programului bazată pe reguli simple și explicabile (pornite din literatura de specialitate, fără
pretenții medicale), un mod de coaching prin care un antrenor poate vedea datele unui client doar cu
acordul lui, și o căutare full-text peste programe și peste istoric. Căutarea este un model de citire
separat, construit peste OpenSearch și ținut la zi prin evenimente declanșate după ce datele se
salvează în baza de date. Partea principală a lucrării este o comparație concretă între OpenSearch și
PostgreSQL, făcută pe un set de date generat și crescut până la 2 milioane de documente. Backend-ul
este scris în Spring Boot~3 / Java~21, cu PostgreSQL~16 ca sursă unică de adevăr, iar corectitudinea
este verificată cu teste de integrare care rulează pe baze de date reale (Testcontainers, fără
mock-uri). Rezultatele arată că cele două soluții se echivalează în jurul a 10.000 de documente și că
cea mai mare diferență apare la căutarea care tolerează greșelile de tastare.
\end{abstract}

\begin{abstract}{english}
This thesis describes a web application for strength training, built around one simple idea: the
training history has to stay correct over time. When a workout is performed, the plan it started from
is copied (``frozen'') as snapshots, so if the user later edits or deletes the data they started
from, what was already recorded does not change. The application lets users manage training programs,
run sessions, follow their progress, use a program marketplace (publish, vote, save, copy), get an
analysis of the program structure based on simple, explainable rules (taken from the training
literature, with no medical claims), use a coaching mode where a coach can see a client's data only
with that client's consent, and search the programs and the history with full-text search. The search
is a separate read model, built on top of OpenSearch and kept up to date through events fired after
the data is saved to the database. The main part of the thesis is a concrete comparison between
OpenSearch and PostgreSQL, run on a generated data set scaled up to 2 million documents. The backend
is written in Spring Boot~3 / Java~21, with PostgreSQL~16 as the single source of truth, and
correctness is checked with integration tests that run on real databases (Testcontainers, no mocks).
The results show that the two options are roughly equivalent around 10{,}000 documents and that the
biggest difference shows up on typo-tolerant (fuzzy) search.
\end{abstract}

\end{abstractpage}
